
\clearpage
\section*{Simulation appendix}
Population evaluation uses independent group quadrature. In orthonormal
Nystr\"om coordinates $\phi(z)$, write
$H_g=\frac13\sum_r\phi(Z_{g,r})\beta(Z_{g,r})'$ and
$\overline H=M^{-1}\sum_g H_g$. Then
\[
\widehat\Sigma=\frac{\sigma_\varepsilon^2}{N}\widehat T_K
+\frac1{GM}\sum_g H_gH_g'
+\left(1-\frac1G\right)\overline H\overline H',
\qquad
\widehat T_K=\frac1{3M}\sum_{g,r}\phi(Z_{g,r})\phi(Z_{g,r})'.
\]
This integrates the factor and idiosyncratic shocks analytically while preserving
within-triplet dependence. The quadrature distribution is independent of all
training paths and cannot affect validation choices. For empirical ranks, full
independent cross-sections are integrated because ranking couples groups.
For heteroskedasticity, the idiosyncratic term uses the actual bounded
state-dependent conditional variance.

\begin{figure}[p]
\centering
\includegraphics[width=\textwidth]{\simoutput/figures/population_loss_complexity.pdf}
\caption{Population ridge loss versus population effective complexity.}
\label{fig:sim_population_loss_complexity}
\end{figure}

\begin{figure}[p]
\centering
\includegraphics[width=\textwidth]{\simoutput/figures/baseline_oos_loss_uncertainty.pdf}
\caption{OOS loss: Monte Carlo mean, median, and pointwise confidence interval for the mean.}
\label{fig:sim_baseline_oos_loss_uncertainty}
\end{figure}

\begin{figure}[p]
\centering
\includegraphics[width=\textwidth]{\simoutput/figures/baseline_oos_sharpe_uncertainty.pdf}
\caption{OOS Sharpe: Monte Carlo mean, median, and pointwise confidence interval for the mean.}
\label{fig:sim_baseline_oos_sharpe_uncertainty}
\end{figure}

\begin{figure}[p]
\centering
\includegraphics[width=\textwidth]{\simoutput/figures/lambda_complexity.pdf}
\caption{Population and empirical effective complexity as reparameterizations of regularization.}
\label{fig:sim_lambda_complexity}
\end{figure}

\begin{figure}[p]
\centering
\includegraphics[width=\textwidth]{\simoutput/figures/selected_complexity_distributions.pdf}
\caption{Selection distributions and oracle versus selected empirical complexity.}
\label{fig:sim_selected_complexity_distributions}
\end{figure}
